\hypersetup{pageanchor=false}
\pagenumbering{gobble}

\begin{titlepage}
\thispagestyle{empty}
\begin{center}
Национальный исследовательский университет \\
«Высшая школа экономики»

\vspace{1cm}

Факультет Программной инженерии

\vfill

{\Large \textbf{Адаптационный курс}}

\vspace{0.5cm}

{\large к дисциплине «Алгоритмы и алгоритмические языки»}

\vspace{1cm}

{\large Методические материалы для студентов первого курса}

\vfill

Москва, 2026
\end{center}
\end{titlepage}

\clearpage
\hypersetup{pageanchor=true}
\pagenumbering{arabic}

\chapter*{Аннотация}
\addcontentsline{toc}{chapter}{Аннотация}

Настоящие методические материалы предназначены для адаптационного курса студентов первого курса направления «Программная инженерия». Курс помогает закрыть пробелы, которые часто остаются после школы, и подготовиться к изучению дисциплины «Алгоритмы и алгоритмические языки», где используются язык C и ассемблер NASM.

В материалах последовательно рассматриваются системы счисления, представление чисел в памяти компьютера, базовая архитектура ЭВМ, операционные системы, Linux как рабочая среда программиста, инструменты сборки и основы системы контроля версий Git.

\chapter*{Как пользоваться материалами}
\addcontentsline{toc}{chapter}{Как пользоваться материалами}

Каждая тема построена вокруг пары занятий: лекции и семинара. Лекция вводит понятия и связывает их с практикой программирования, а семинар закрепляет материал через задачи, команды и небольшие эксперименты.

Во второй главе темы дополнительно разделены на обязательные и усложнённые. Обязательная часть задаёт минимальный уровень подготовки для дальнейшего изучения C и NASM. Усложнённая часть предназначена для студентов, которым нужна большая глубина или дополнительные ориентиры для самостоятельной работы.

\tableofcontents
